\section{Introduction}

RECIST v1.1~\cite{eisenhauerNewResponseEvaluation2009} is widely used to standardize tumor response assessment in clinical practice and clinical trials. Its diameter-based measurements are practical and easy to acquire, but a small number of 2D measurements cannot fully represent the 3D shape and spatial extent of a lesion. Volumetric lesion segmentation can describe tumor burden more completely, and previous studies suggest that volumetric analysis may provide clinically useful information beyond RECIST in settings such as lung cancer response assessment and breast cancer treatment response prediction~\cite{hayesComparisonCTVolumetric2016,machadoPromptableCTFoundation2025}. However, full 3D lesion annotation is difficult to scale in time-constrained clinical workflows.

The CVPR 2026: Foundation Models for Pan-cancer Segmentation in CT Images competition addresses this gap by asking models to predict a 3D lesion mask from a CT volume and a RECIST marker. This setting requires accurate local segmentation from sparse spatial input, while lesions vary in size, shape, appearance, and anatomical location. The challenge also emphasizes efficient inference under strict computational constraints, making runtime and memory use important design considerations.

\subsection{Related work}

Medical image segmentation has long been dominated by encoder--decoder architectures, most notably U-Net~\cite{ronnebergerUNetConvolutionalNetworks2015}. 3D variants such as SegResNet~\cite{myronenko3DMRIBrain2018} extend this family to volumetric segmentation, while nnU-Net~\cite{isenseeNnUNetSelfadaptingFramework2018} showed the importance of adapting preprocessing, training, inference, and postprocessing to each dataset. U-Net-style models therefore remain strong baselines for medical segmentation, especially when combined with task-specific configuration~\cite{maFastLowResourceAccurate2026}.

Recent ConvNeXt-style models further show that convolutional backbones remain competitive for 3D medical image segmentation. MedNeXt~\cite{royMedNeXtTransformerDrivenScaling2023} adapts ConvNeXt~\cite{liuConvNet2020s2022} principles to volumetric segmentation through a fully convolutional encoder--decoder with residual upsampling and downsampling blocks. MedNeXt-v2~\cite{royMedNeXtv2Scaling3D2025} extends this line of work by scaling 3D ConvNeXt backbones for supervised representation learning and adding volumetric global response normalization. Its evaluation against several strong public segmentation models~\cite{wasserthalTotalSegmentatorRobustSegmentation2023,hantzeMRSegmentatorMultiModalitySegmentation2025,heVISTA3DUnifiedSegmentation2024,duSegVolUniversalInteractive2025,huangSTUNetScalableTransferable2023} makes it a useful reference point for modern 3D medical segmentation backbones.

Promptable segmentation conditions prediction models on user- or task-provided spatial information. General segmentation models established this paradigm for natural images, videos, and concept-level prompts~\cite{kirillovSegmentAnything2023,raviSAM2Segment2024,carionSAM3Segment2026}, while medical variants adapt promptable segmentation to medical images and volumetric data~\cite{maSegmentAnythingMedical2024,wangSAMMed3DGeneralpurposeSegmentation2024,duSegVolUniversalInteractive2025}. In 3D medical segmentation, nnInteractive is especially relevant because it represents spatial prompts as additional input channels and introduces AutoZoom for adaptive field-of-view refinement~\cite{isenseeNnInteractiveRedefining3D2025}.

ENSAM introduced a promptable 3D medical image segmentation framework with a SegResNet-based image encoder, prompt encoder, and mask decoder connected through latent cross-attention~\cite{stenhedeENSAMEfficientFoundation2025}. Lite ENSAM adapted this framework to RECIST-conditioned CT lesion segmentation by using RECIST markers as sparse prompts and by reducing the computational footprint for CPU-based inference~\cite{bjornstadLiteENSAMLightweight2025}. Together, these models provide the main starting point for the present work: prompt-conditioned volumetric segmentation, RECIST-guided lesion localization, and lightweight inference. 

Building on ENSAM and Lite ENSAM, we present LETT-NeXt, a lightweight RECIST-guided 3D lesion segmentation model. LETT-NeXt replaces latent prompt decoding with direct RECIST prompt-channel conditioning and uses a compact MedNeXt-v2 encoder–decoder for efficient local segmentation. During training, an auxiliary anatomy–tumor head provides additional supervision. At inference, RECIST-guided regional crops and AutoZoom maintain a local field of view, while still allowing adaptive expansion when needed.



% The incremental additions of this work are:
% \begin{enumerate}
%     \item Backbone improvement: the previous SegResNet backbone is replaced with a NeXt backbone, improving \textcolor{red}{XXX}.
%     \item Architectural simplification: The latent cross-attention pattern from ENSAM and Lite-ENSAM is replaced with a dense \textcolor{red}{Describe?}
%     \item PanTS auxiliary supervision: uses auxiliary anatomy–tumor supervision during training to improve lesion segmentation.
% \end{enumerate}